\chapter{Cơ sở lý thuyết}

Chương này trình bày các nền tảng lý thuyết được sử dụng trong đề tài, bao gồm: bài toán Quy hoạch mở rộng nguồn điện (GEP), tối ưu hóa hai tầng (Bi-level Optimization), quy hoạch tuyến tính (LP), tối ưu hóa Bayes (Bayesian Optimization), và học tăng cường (Reinforcement Learning). Mỗi phần lý thuyết đều được phân tích trong mối liên hệ trực tiếp với bài toán đang giải quyết.

% ================================================================
\section{Bài toán Quy hoạch mở rộng nguồn điện (GEP)}
% ================================================================

\subsection{Định nghĩa bài toán}

Bài toán Quy hoạch mở rộng nguồn điện (Generation Expansion Planning -- GEP) là bài toán xác định thời điểm, vị trí, loại công nghệ và quy mô công suất của các nhà máy phát điện mới cần xây dựng trong một khoảng thời gian quy hoạch nhất định, sao cho tổng chi phí hệ thống (bao gồm chi phí đầu tư và chi phí vận hành) được tối thiểu hóa, đồng thời thỏa mãn các ràng buộc về kỹ thuật, kinh tế và môi trường \cite{koltsaklis2018gep}.

Về mặt toán học, bài toán GEP tổng quát có thể được phát biểu như sau:

\begin{equation}
    \min_{\mathbf{x}} \quad C_{\text{inv}}(\mathbf{x}) + C_{\text{op}}(\mathbf{x})
    \label{eq:gep_general}
\end{equation}

\noindent thỏa mãn các ràng buộc:
\begin{itemize}
    \item Cân bằng cung -- cầu công suất tại mọi thời điểm
    \item Giới hạn công suất phát của từng nhà máy
    \item Giới hạn truyền tải giữa các vùng
    \item Ràng buộc về độ tin cậy cung cấp điện
    \item Ràng buộc về phát thải (nếu có)
\end{itemize}

\subsection{Phân loại các phương pháp giải GEP}

Theo tổng quan của Sadeghi và cộng sự \cite{sadeghi2017gep_review}, các phương pháp giải GEP được chia thành ba nhóm chính:

\begin{enumerate}
    \item \textbf{Phương pháp toán học chính xác:} Bao gồm Quy hoạch tuyến tính (LP), Quy hoạch nguyên hỗn hợp (MILP), và Quy hoạch động (DP). Các phương pháp này đảm bảo tìm được nghiệm tối ưu toàn cục nhưng gặp khó khăn khi quy mô bài toán tăng lên (hiện tượng ``lời nguyền chiều'' -- curse of dimensionality) \cite{oree2017gep_review}.
    
    \item \textbf{Phương pháp metaheuristic:} Bao gồm Thuật toán di truyền (GA), Tối ưu bầy đàn (PSO), Mô phỏng luyện kim (SA). Các phương pháp này linh hoạt hơn nhưng không đảm bảo hội tụ về nghiệm tối ưu toàn cục.
    
    \item \textbf{Phương pháp dựa trên AI:} Bao gồm Mạng thần kinh nhân tạo, Học tăng cường, và các phương pháp lai. Đây là hướng nghiên cứu mới và được quan tâm ngày càng nhiều nhờ khả năng xử lý các bài toán quy mô lớn với tính bất định cao.
\end{enumerate}

\iffalse
\subsection{Liên hệ với đề tài}

Trong đề tài này, bài toán GEP được giải quyết bằng cách kết hợp hai phương pháp: tầng trên sử dụng các thuật toán AI (Bayesian Optimization và RL) để khám phá không gian quyết định đầu tư, trong khi tầng dưới sử dụng LP để đảm bảo tính khả thi vật lý. Cách tiếp cận lai này nhằm khắc phục nhược điểm của từng phương pháp khi được sử dụng độc lập: LP thuần túy không thể xử lý hiệu quả không gian quyết định đầu tư rộng lớn, trong khi AI thuần túy khó đảm bảo các ràng buộc vật lý nghiêm ngặt.
\fi

% ================================================================
\section{Tối ưu hóa hai tầng (Bi-level Optimization)}
% ================================================================

\subsection{Khái niệm và phát biểu toán học}

Tối ưu hóa hai tầng (Bi-level Optimization) là một lớp bài toán tối ưu hóa toán học trong đó một bài toán tối ưu (tầng trên) chứa một bài toán tối ưu khác (tầng dưới) như một ràng buộc \cite{colson2007bilevel}. Dạng tổng quát của bài toán bi-level được phát biểu như sau:

\begin{align}
    &\min_{\mathbf{x}, \mathbf{y}} \quad F(\mathbf{x}, \mathbf{y}) \quad &\text{(Hàm mục tiêu tầng trên)} \label{eq:bilevel_upper}\\
    &\text{s.t.} \quad G(\mathbf{x}, \mathbf{y}) \leq 0 \quad &\text{(Ràng buộc tầng trên)} \notag \\
    &\quad\quad\;\; \mathbf{y} \in \arg\min_{\mathbf{y}'} \left\{ f(\mathbf{x}, \mathbf{y}') \;:\; g(\mathbf{x}, \mathbf{y}') \leq 0 \right\} \quad &\text{(Bài toán tầng dưới)} \label{eq:bilevel_lower}
\end{align}

\noindent trong đó $\mathbf{x}$ là biến quyết định của tầng trên, $\mathbf{y}$ là biến quyết định của tầng dưới phụ thuộc vào lựa chọn $\mathbf{x}$ của tầng trên.

\subsection{Đặc điểm của bài toán bi-level}

Bài toán bi-level có một số đặc điểm quan trọng \cite{sinha2018bilevel_review}:

\begin{itemize}
    \item \textbf{Cấu trúc phân cấp:} Hai bài toán tối ưu có mối quan hệ chủ -- phụ (leader -- follower). Tầng trên đưa ra quyết định trước, tầng dưới phản hồi tối ưu dựa trên quyết định của tầng trên.
    
    \item \textbf{Tính NP-hard:} Ngay cả khi cả hai tầng đều là bài toán tuyến tính, bài toán bi-level tổng quát vẫn thuộc lớp NP-hard \cite{hansen1992bilevel_nphard}.
    
    \item \textbf{Phản hồi ngầm:} Nghiệm tối ưu $\mathbf{y}^*(\mathbf{x})$ của tầng dưới là một hàm ngầm (implicit function) của $\mathbf{x}$, khiến hàm mục tiêu tầng trên thường không khả vi (non-differentiable) hoặc phi lồi (non-convex).
\end{itemize}

\iffalse
\subsection{Liên hệ với đề tài}

Trong đề tài này, cấu trúc bi-level được hiện thực hóa như sau:

\begin{itemize}
    \item \textbf{Tầng trên (Leader):} Quyết định véctơ đầu tư $\mathbf{x} = [x_{z,k}]$ -- tức công suất MW thêm mới cho mỗi công nghệ $k$ tại mỗi vùng $z$. Hàm mục tiêu tầng trên là tối thiểu hóa tổng chi phí (đầu tư + vận hành + phạt cắt tải).
    
    \item \textbf{Tầng dưới (Follower):} Với mỗi quyết định đầu tư $\mathbf{x}$ từ tầng trên, tầng dưới giải bài toán điều độ kinh tế (Economic Dispatch) bằng LP để xác định phương thức vận hành tối ưu, bao gồm phân bổ công suất phát, truyền tải liên vùng, và cắt tải (nếu cần thiết).
\end{itemize}

Mối quan hệ phản hồi giữa hai tầng tạo thành vòng lặp: tầng trên đề xuất $\mathbf{x}$, tầng dưới trả về chi phí vận hành $C_{\text{op}}(\mathbf{x})$ và lượng cắt tải $S_{\text{total}}(\mathbf{x})$, tầng trên sử dụng thông tin này để cải thiện quyết định đầu tư ở lần lặp tiếp theo.
\fi

% ================================================================
\section{Quy hoạch tuyến tính (Linear Programming)}
% ================================================================

\subsection{Lý thuyết cơ bản}

Quy hoạch tuyến tính (Linear Programming -- LP) là phương pháp tối ưu hóa một hàm mục tiêu tuyến tính trên một miền chấp nhận được xác định bởi các ràng buộc tuyến tính \cite{bertsimas1997lp}. Dạng chuẩn (standard form) của bài toán LP:

\begin{align}
    \min_{\mathbf{y}} \quad & \mathbf{c}^\top \mathbf{y} \label{eq:lp_obj}\\
    \text{s.t.} \quad & A\mathbf{y} = \mathbf{b} \notag\\
    & \mathbf{y} \geq \mathbf{0} \notag
\end{align}

\noindent trong đó $\mathbf{c} \in \mathbb{R}^n$ là véctơ chi phí, $A \in \mathbb{R}^{m \times n}$ là ma trận ràng buộc, $\mathbf{b} \in \mathbb{R}^m$ là véctơ vế phải.

\subsection{Thuật toán giải LP}

Hai thuật toán chính để giải LP là:

\begin{enumerate}
    \item \textbf{Phương pháp Đơn hình (Simplex Method):} Do Dantzig đề xuất năm 1947 \cite{dantzig1951simplex}, phương pháp này duyệt qua các đỉnh của đa diện chấp nhận được. Mặc dù có độ phức tạp trường hợp xấu nhất là hàm mũ, trên thực tế Simplex thường hoạt động rất hiệu quả.
    
    \item \textbf{Phương pháp Điểm trong (Interior Point Method):} Do Karmarkar đề xuất năm 1984 \cite{karmarkar1984interior}, phương pháp này di chuyển qua phần bên trong của miền chấp nhận được. Có độ phức tạp đa thức $O(n^{3.5}L)$ với $L$ là kích thước đầu vào.
\end{enumerate}

\subsection{Bài toán Điều độ kinh tế (Economic Dispatch)}

Bài toán Điều độ kinh tế (Economic Dispatch -- ED) là bài toán phân bổ công suất phát cho các nhà máy phát điện sao cho tổng chi phí vận hành là nhỏ nhất, đồng thời thỏa mãn cân bằng cung -- cầu và các ràng buộc kỹ thuật \cite{wood2013powersystems}. Khi hàm chi phí được xấp xỉ tuyến tính, ED trở thành một bài toán LP.

Trong đề tài, bài toán ED theo từng giờ $t$ được phát biểu:

\begin{equation}
    \min_{P_i^t,\; F_{zz'}^t,\; \Delta L_z^t} \quad
    \sum_{i \in \mathbf{G}} c_i \, P_i^t + \beta \sum_{z \in \mathbf{Z}} \Delta L_z^t
    \label{eq:ed_objective}
\end{equation}

\noindent thỏa mãn:

\noindent \textit{(i) Cân bằng công suất từng vùng:}
\begin{equation}
    \sum_{i \in \mathbf{G}_z} P_i^t + \sum_{z' \ne z} F_{z'z}^t - \sum_{z' \ne z} F_{zz'}^t + \Delta L_z^t = D_z^t,
    \quad \forall z \in \mathbf{Z}
    \label{eq:power_balance}
\end{equation}

\noindent \textit{(ii) Giới hạn công suất phát:}
\begin{equation}
    0 \leq P_i^t \leq P_i^{\max} \cdot \text{CF}_i(t),
    \quad \forall i \in \mathbf{G}
    \label{eq:gen_limit}
\end{equation}

\noindent \textit{(iii) Giới hạn truyền tải liên vùng:}
\begin{equation}
    -F_{zz'}^{\max} \leq F_{zz'}^t \leq F_{zz'}^{\max},
    \quad \forall (z, z') \in \mathbf{L}
    \label{eq:transfer_limit}
\end{equation}

\noindent \textit{(iv) Biến cắt tải không âm:}
\begin{equation}
    \Delta L_z^t \geq 0,
    \quad \forall z \in \mathbf{Z}
    \label{eq:shed_nonneg}
\end{equation}

\noindent trong đó:
\begin{itemize}
    \item $P_i^t$: Công suất phát của tổ máy $i$ tại giờ $t$ (MW)
    \item $c_i$: Chi phí biên vận hành của tổ máy $i$ (\$/MWh)
    \item $F_{zz'}^t$: Công suất truyền tải từ vùng $z$ sang vùng $z'$ tại giờ $t$ (MW)
    \item $\Delta L_z^t$: Lượng cắt tải tại vùng $z$, giờ $t$ (MW)
    \item $D_z^t$: Phụ tải tại vùng $z$, giờ $t$ (MW)
    \item $\text{CF}_i(t) \in [0, 1]$: Hệ số công suất (Capacity Factor) theo giờ
    \item $\beta = \text{VOLL}$: Value of Lost Load -- giá trị kinh tế của mỗi MWh bị cắt tải
\end{itemize}

\iffalse
\subsection{Liên hệ với đề tài}

LP đóng vai trò là ``bộ mô phỏng'' ở tầng dưới của hệ thống bi-level. Mỗi khi tầng trên đề xuất một phương án đầu tư $\mathbf{x}$, tầng dưới:
\begin{enumerate}
    \item Cập nhật danh sách nguồn phát: thêm các nhà máy mới với công suất $x_{z,k}$ vào từng vùng $z$
    \item Giải bài toán LP (\ref{eq:ed_objective}) cho tất cả $T$ giờ mô phỏng
    \item Tổng hợp chi phí vận hành $C_{\text{op}}(\mathbf{x})$ và lượng cắt tải $S_{\text{total}}(\mathbf{x})$
    \item Trả kết quả về tầng trên
\end{enumerate}

LP được chọn nhờ hai ưu điểm quan trọng: \textit{(i)} luôn tìm được nghiệm tối ưu toàn cục (nếu bài toán khả thi), và \textit{(ii)} tốc độ giải rất nhanh nhờ các solver hiện đại như HiGHS \cite{huangfu2018highs}, phù hợp cho việc đánh giá lặp đi lặp lại trong vòng lặp bi-level.
\fi

% ================================================================
\section{Tối ưu hóa Bayes (Bayesian Optimization)}
% ================================================================

\subsection{Khái niệm và động lực}

Tối ưu hóa Bayes (Bayesian Optimization -- BO) là phương pháp tối ưu hóa toàn cục dành cho các hàm mục tiêu dạng \textit{hộp đen} (black-box function) có chi phí đánh giá đắt đỏ \cite{shahriari2016bo_survey}. Khác với các phương pháp dựa trên gradient, BO không yêu cầu biết dạng giải tích hay đạo hàm của hàm mục tiêu, mà thay vào đó xây dựng một mô hình xác suất xấp xỉ (surrogate model) dựa trên các quan sát đã thu thập.

Bài toán BO được phát biểu:
\begin{equation}
    \mathbf{x}^* = \arg\min_{\mathbf{x} \in \mathbf{X}} f(\mathbf{x})
    \label{eq:bo_problem}
\end{equation}
\noindent trong đó $f$ là hàm mục tiêu đắt đỏ (expensive-to-evaluate), $\mathbf{X}$ là miền tìm kiếm, và ta không có quyền truy cập vào gradient $\nabla f$.

\subsection{Các thành phần cốt lõi}

BO gồm hai thành phần chính \cite{frazier2018bo_tutorial}:

\subsubsection{Mô hình thay thế (Surrogate Model)}

Mô hình thay thế là một mô hình xác suất xấp xỉ hàm mục tiêu $f(\mathbf{x})$ dựa trên tập dữ liệu đã quan sát $\mathbf{D}_n = \{(\mathbf{x}_i, y_i)\}_{i=1}^n$. Hai lựa chọn phổ biến nhất:

\begin{itemize}
    \item \textbf{Gaussian Process (GP):} Mô hình GP định nghĩa một phân phối trên không gian hàm, nơi mỗi tập con hữu hạn tuân theo phân phối chuẩn đa biến \cite{rasmussen2006gp}:
    \begin{equation}
        f(\mathbf{x}) \sim \mathrm{GP}\left(m(\mathbf{x}),\; k(\mathbf{x}, \mathbf{x}')\right)
    \end{equation}
    với $m(\mathbf{x})$ là hàm trung bình và $k(\mathbf{x}, \mathbf{x}')$ là hàm hiệp phương sai (kernel).
    
    \item \textbf{Tree-structured Parzen Estimator (TPE):} Do Bergstra và cộng sự đề xuất \cite{bergstra2011tpe}, TPE mô hình hóa hai mật độ xác suất có điều kiện:
    \begin{equation}
        p(\mathbf{x} \mid y) = \begin{cases}
            l(\mathbf{x}) & \text{nếu } y < y^* \\
            g(\mathbf{x}) & \text{nếu } y \geq y^*
        \end{cases}
        \label{eq:tpe}
    \end{equation}
    trong đó $y^*$ là ngưỡng phân chia (thường là quantile), $l(\mathbf{x})$ là mật độ các điểm ``tốt'', $g(\mathbf{x})$ là mật độ các điểm ``kém''. Tỷ số $l(\mathbf{x})/g(\mathbf{x})$ cao tương ứng với vùng có triển vọng tốt.
\end{itemize}

\subsubsection{Hàm thu thập (Acquisition Function)}

Hàm thu thập $\alpha(\mathbf{x})$ quyết định điểm nào sẽ được đánh giá tiếp theo, bằng cách cân bằng giữa \textit{khai thác} (exploitation) -- thử nghiệm tại vùng đã biết tốt, và \textit{khám phá} (exploration) -- thử nghiệm tại vùng chưa biết nhiều.

\textbf{Expected Improvement (EI)} là hàm thu thập phổ biến nhất \cite{jones1998ei}:
\begin{equation}
    \text{EI}(\mathbf{x}) = \mathbb{E}\left[\max\left(f_{\min} - f(\mathbf{x}),\; 0\right)\right]
    \label{eq:ei}
\end{equation}
\noindent trong đó $f_{\min}$ là giá trị tốt nhất đã quan sát được. Điểm tiếp theo được chọn bằng:
\begin{equation}
    \mathbf{x}_{n+1} = \arg\max_{\mathbf{x} \in \mathbf{X}} \text{EI}(\mathbf{x})
\end{equation}

\subsection{Quy trình tối ưu hóa Bayes}

Thuật toán BO hoạt động theo các bước vòng lặp tuần tự cơ bản sau:

\begin{enumerate}
    \item Khởi tạo: Đánh giá hàm $f$ tại $n_0$ điểm ngẫu nhiên trong $\mathbf{X}$
    \item Lặp cho đến khi hết ngân sách (budget):
    \begin{enumerate}
        \item Cập nhật mô hình thay thế dựa trên $\mathbf{D}_n$
        \item Tối ưu hàm thu thập $\alpha(\mathbf{x})$ để tìm $\mathbf{x}_{n+1}$
        \item Đánh giá $y_{n+1} = f(\mathbf{x}_{n+1})$
        \item Cập nhật $\mathbf{D}_{n+1} = \mathbf{D}_n \cup \{(\mathbf{x}_{n+1}, y_{n+1})\}$
    \end{enumerate}
    \item Trả về $\mathbf{x}^* = \arg\min_{(\mathbf{x}_i, y_i) \in \mathbf{D}} y_i$
\end{enumerate}

\iffalse
\subsection{Liên hệ với đề tài}

Trong đề tài, BO được sử dụng ở tầng trên để tìm véctơ đầu tư $\mathbf{x}^*$ tối ưu. Hàm mục tiêu ``hộp đen'' chính là điểm phạt $J(\mathbf{x})$ từ tầng dưới -- mỗi lần đánh giá yêu cầu giải bài toán LP cho hàng ngàn giờ mô phỏng. Cụ thể:

\begin{itemize}
    \item \textbf{Mô hình thay thế:} Sử dụng TPE (phương trình \ref{eq:tpe}) thông qua thư viện Optuna \cite{akiba2019optuna}, phù hợp hơn GP cho không gian tham số hỗn hợp (liên tục + rời rạc) và số chiều trung bình.
    
    \item \textbf{Không gian tìm kiếm:} $\mathbf{X} = \prod_{z \in \mathbf{Z}} \prod_{k \in \mathbf{K}} [0, \overline{X}_k]$ với bước rời rạc hóa 10 MW để giảm kích thước hiệu dụng.
    
    \item \textbf{Hiệu quả mẫu:} BO chỉ cần 100--200 lần đánh giá để hội tụ, so với hàng ngàn episodes cần thiết cho RL, giúp tiết kiệm đáng kể thời gian tính toán.
\end{itemize}
\fi

% ================================================================
\section{Quá trình quyết định Markov (MDP)}
% ================================================================

\subsection{Định nghĩa}

Quá trình quyết định Markov (Markov Decision Process -- MDP) là khung toán học mô tả bài toán ra quyết định tuần tự trong môi trường có tính ngẫu nhiên \cite{puterman1994mdp}. Một MDP được xác định bởi bộ năm $(\mathbf{S}, \mathbf{A}, P, R, \gamma)$:

\begin{itemize}
    \item $\mathbf{S}$: Không gian trạng thái (state space)
    \item $\mathbf{A}$: Không gian hành động (action space)
    \item $P(s' \mid s, a)$: Hàm chuyển trạng thái (transition probability)
    \item $R(s, a)$: Hàm phần thưởng (reward function)
    \item $\gamma \in [0, 1]$: Hệ số chiết khấu (discount factor)
\end{itemize}

\noindent \textbf{Tính chất Markov} cho rằng trạng thái tương lai chỉ phụ thuộc vào trạng thái hiện tại, không phụ thuộc vào lịch sử:
\begin{equation}
    P(s_{t+1} \mid s_t, a_t, s_{t-1}, a_{t-1}, \ldots) = P(s_{t+1} \mid s_t, a_t)
\end{equation}

\subsection{Chính sách và hàm giá trị}

\textbf{Chính sách (Policy)} $\pi: \mathbf{S} \to \mathbf{P}(\mathbf{A})$ ánh xạ mỗi trạng thái sang một phân phối xác suất trên không gian hành động. Mục tiêu là tìm chính sách $\pi^*$ cực đại hóa tổng phần thưởng kỳ vọng:

\begin{equation}
    \pi^* = \arg\max_{\pi} \; \mathbb{E}_{\pi}\left[\sum_{t=0}^{\infty} \gamma^t R(s_t, a_t)\right]
    \label{eq:optimal_policy}
\end{equation}

\textbf{Hàm giá trị trạng thái (State-value function):}
\begin{equation}
    V^{\pi}(s) = \mathbb{E}_{\pi}\left[\sum_{t=0}^{\infty} \gamma^t R(s_t, a_t) \;\middle|\; s_0 = s\right]
    \label{eq:value_function}
\end{equation}

\textbf{Hàm giá trị hành động (Action-value function):}
\begin{equation}
    Q^{\pi}(s, a) = \mathbb{E}_{\pi}\left[\sum_{t=0}^{\infty} \gamma^t R(s_t, a_t) \;\middle|\; s_0 = s, a_0 = a\right]
    \label{eq:q_function}
\end{equation}

\textbf{Phương trình Bellman:} Mối quan hệ đệ quy giữa giá trị của các trạng thái kế tiếp:
\begin{equation}
    V^{\pi}(s) = \sum_{a} \pi(a|s) \left[ R(s,a) + \gamma \sum_{s'} P(s'|s,a) V^{\pi}(s') \right]
    \label{eq:bellman}
\end{equation}

\iffalse
\subsection{Liên hệ với đề tài}

Bài toán GEP được ánh xạ thành MDP đơn bước (single-step MDP, hay còn gọi là bandit contextual) như sau:

\begin{itemize}
    \item \textbf{Trạng thái $\mathbf{s}$:} Véctơ đặc trưng hệ thống (công suất hiện hữu, phụ tải đỉnh/trung bình, biên dự phòng) đã chuẩn hóa về $[-1, 1]$.
    \item \textbf{Hành động $\mathbf{a}$:} Véctơ đầu tư $\mathbf{x} = [x_{z,k}]$ -- công suất MW thêm mới.
    \item \textbf{Phần thưởng $r$:} $r = -J(\mathbf{x})/10^6$ -- nghịch đảo điểm phạt từ tầng dưới.
\end{itemize}

Việc xây dựng khung MDP cho phép áp dụng trực tiếp các thuật toán RL hiện đại (PPO, SAC) mà không cần thiết kế heuristic riêng, đồng thời mở đường cho mở rộng sang quy hoạch đa giai đoạn trong tương lai.
\fi

% ================================================================
\section{Học tăng cường (Reinforcement Learning)}
% ================================================================

\subsection{Tổng quan}

Học tăng cường (Reinforcement Learning -- RL) là nhánh của machine learning trong đó một tác tử (agent) học cách hành động thông qua tương tác trực tiếp với môi trường, nhận tín hiệu phần thưởng để cải thiện chính sách theo thời gian \cite{sutton2018rl}. Khác với học có giám sát (supervised learning), RL không yêu cầu dữ liệu huấn luyện với nhãn cho trước, mà tác tử tự khám phá và khai thác.

\subsection{Phương pháp Policy Gradient}

Trong các bài toán với không gian hành động liên tục (như bài toán GEP), các phương pháp Policy Gradient được ưu tiên sử dụng \cite{schulman2015trpo}. Ý tưởng cốt lõi là tham số hóa chính sách $\pi_\theta(a|s)$ bằng mạng thần kinh với tham số $\theta$, và cập nhật $\theta$ theo hướng gradient của hàm mục tiêu:

\begin{equation}
    \nabla_\theta J(\theta) = \mathbb{E}_{\pi_\theta}\left[\nabla_\theta \log \pi_\theta(a_t | s_t) \cdot A^{\pi_\theta}(s_t, a_t)\right]
    \label{eq:policy_gradient}
\end{equation}

\noindent trong đó $A^{\pi}(s, a) = Q^{\pi}(s, a) - V^{\pi}(s)$ là hàm lợi thế (advantage function), đo lường mức độ ``tốt hơn trung bình'' của hành động $a$ tại trạng thái $s$.

\subsection{PPO -- Proximal Policy Optimization}

PPO do Schulman và cộng sự đề xuất năm 2017 \cite{schulman2017ppo} là một thuật toán policy gradient với cơ chế giới hạn bước cập nhật chính sách, đảm bảo quá trình huấn luyện ổn định.

\subsubsection{Hàm mục tiêu có kẹp (Clipped Objective)}

PPO sử dụng hàm mục tiêu có kẹp (clipped surrogate objective):

\begin{equation}
    L^{\text{CLIP}}(\theta) = \hat{\mathbb{E}}_t \left[ \min\left( r_t(\theta) \hat{A}_t,\; \text{clip}\left(r_t(\theta), 1-\epsilon, 1+\epsilon\right) \hat{A}_t \right) \right]
    \label{eq:ppo_clip}
\end{equation}

\noindent trong đó:
\begin{itemize}
    \item $r_t(\theta) = \frac{\pi_\theta(a_t | s_t)}{\pi_{\theta_{\text{old}}}(a_t | s_t)}$ là tỷ số xác suất (probability ratio)
    \item $\hat{A}_t$ là ước lượng hàm lợi thế (thường dùng GAE -- Generalized Advantage Estimation \cite{schulman2015gae})
    \item $\epsilon$ là siêu tham số kẹp (thường $\epsilon = 0.2$)
\end{itemize}

Phép kẹp ngăn chặn các bước cập nhật quá lớn: khi $r_t(\theta)$ ra ngoài khoảng $[1-\epsilon, 1+\epsilon]$, gradient bị triệt tiêu, bảo vệ quá trình huấn luyện khỏi sự bất ổn.

\subsubsection{Ưu điểm của PPO}
\begin{itemize}
    \item Đơn giản trong triển khai so với TRPO (Trust Region Policy Optimization)
    \item Ổn định và ít nhạy cảm với siêu tham số
    \item Hiệu quả trên nhiều loại bài toán khác nhau
    \item Thuộc loại \textit{on-policy}: sử dụng dữ liệu mới nhất từ chính sách hiện tại
\end{itemize}

\subsection{SAC -- Soft Actor-Critic}

SAC do Haarnoja và cộng sự đề xuất năm 2018 \cite{haarnoja2018sac} là thuật toán RL off-policy dựa trên nguyên lý cực đại hóa entropy (maximum entropy RL).

\subsubsection{Nguyên lý cực đại hóa Entropy}

SAC tối ưu hàm mục tiêu mở rộng bao gồm cả entropy của chính sách:

\begin{equation}
    \pi^* = \arg\max_\pi \; \mathbb{E}_\pi \left[ \sum_{t=0}^{\infty} \gamma^t \left( R(s_t, a_t) + \alpha H\left(\pi(\cdot | s_t)\right) \right) \right]
    \label{eq:sac_objective}
\end{equation}

\noindent trong đó $H(\pi(\cdot|s)) = -\mathbb{E}_{a \sim \pi}[\log \pi(a|s)]$ là entropy Shannon của chính sách, và $\alpha > 0$ là hệ số nhiệt độ (temperature coefficient) điều chỉnh mức độ khám phá.

\subsubsection{Kiến trúc Actor-Critic}

SAC sử dụng kiến trúc Actor-Critic với ba thành phần mạng neuron:

\begin{itemize}
    \item \textbf{Actor $\pi_\phi$:} Xuất ra phân phối hành động (thường là Gaussian) dựa trên trạng thái. Hàm cập nhật:
    \begin{equation}
        J_\pi(\phi) = \mathbb{E}_{s_t \sim \mathbf{D}} \left[ \mathbb{E}_{a_t \sim \pi_\phi} \left[ \alpha \log \pi_\phi(a_t | s_t) - Q_\psi(s_t, a_t) \right] \right]
    \end{equation}
    
    \item \textbf{Critic $Q_\psi$ (hai mạng):} Ước lượng hàm Q-value mềm (soft Q-function). Sử dụng hai mạng critic ($Q_{\psi_1}$, $Q_{\psi_2}$) và lấy min để giảm overestimation bias:
    \begin{equation}
        J_Q(\psi_i) = \mathbb{E}_{(s_t, a_t) \sim \mathbf{D}} \left[ \frac{1}{2} \left( Q_{\psi_i}(s_t, a_t) - \hat{Q}(s_t, a_t) \right)^2 \right]
    \end{equation}
    với target value:
    \begin{equation}
        \hat{Q}(s_t, a_t) = R(s_t, a_t) + \gamma \mathbb{E}_{s_{t+1}} \left[ \min_{j=1,2} Q_{\bar{\psi}_j}(s_{t+1}, \tilde{a}_{t+1}) - \alpha \log \pi_\phi(\tilde{a}_{t+1} | s_{t+1}) \right]
    \end{equation}
\end{itemize}

\subsubsection{Ưu điểm của SAC so với PPO}
\begin{itemize}
    \item \textbf{Off-policy:} Có thể tái sử dụng dữ liệu từ replay buffer, tăng hiệu quả mẫu (sample efficiency)
    \item \textbf{Khám phá tự động:} Entropy regularization giúp chính sách duy trì tính đa dạng, tránh hội tụ sớm về cực trị địa phương
    \item \textbf{Tự điều chỉnh $\alpha$:} Phiên bản cải tiến tự động điều chỉnh hệ số nhiệt độ $\alpha$ dựa trên entropy mục tiêu \cite{haarnoja2018sac_v2}
\end{itemize}

\iffalse
\subsection{Liên hệ với đề tài}

Trong đề tài, cả PPO và SAC đều được triển khai ở tầng trên thông qua thư viện Stable-Baselines3 \cite{raffin2021sb3}. Mỗi thuật toán đóng vai trò như một ``nhà lập kế hoạch'' thông minh:

\begin{itemize}
    \item PPO phù hợp cho các kịch bản đơn giản, cần sự ổn định trong huấn luyện
    \item SAC phù hợp cho các kịch bản phức tạp, cần khám phá rộng không gian đầu tư nhờ cơ chế entropy
    \item Cả hai đều tương tác với tầng dưới (LP) qua giao diện OpenAI Gymnasium \cite{towers2024gymnasium}: xuất hành động (véctơ đầu tư), nhận phần thưởng (nghịch đảo điểm phạt)
\end{itemize}
\fi

% ================================================================
\section{Chi phí đầu tư niên kim hóa (Annualized CAPEX)}
% ================================================================

\subsection{Nhân tử thu hồi vốn (Capital Recovery Factor)}

Để so sánh công bằng giữa chi phí đầu tư một lần và chi phí vận hành hàng năm, chi phí đầu tư được quy đổi thành dòng tiền hàng năm đều nhau trong suốt vòng đời kinh tế của công nghệ. Công cụ toán học sử dụng là Nhân tử thu hồi vốn (Capital Recovery Factor -- CRF) \cite{brealey2020corporate}:

\begin{equation}
    \text{CRF}(r, n) = \frac{r(1+r)^n}{(1+r)^n - 1}
    \label{eq:crf}
\end{equation}

\noindent trong đó:
\begin{itemize}
    \item $r$: Lãi suất chiết khấu (WACC -- Weighted Average Cost of Capital), trong đề tài $r = 7\%$
    \item $n$: Tuổi thọ kinh tế của công nghệ (năm)
\end{itemize}

Chi phí đầu tư niên kim hóa:
\begin{equation}
    C_{\text{inv}}(\mathbf{x}) = \sum_{z \in \mathbf{Z}} \sum_{k \in \mathbf{K}} x_{z,k} \cdot \text{CAPEX}_k \cdot \text{CRF}(r, n_k)
    \label{eq:annualized_capex}
\end{equation}

\iffalse
\subsection{Liên hệ với đề tài}

CRF cho phép quy đổi CAPEX (chi phí đầu tư ban đầu) thành dòng chi phí hàng năm, từ đó có thể cộng trực tiếp với chi phí vận hành hàng năm để tạo thành hàm mục tiêu thống nhất. Điều này đặc biệt quan trọng khi so sánh các công nghệ có tuổi thọ khác nhau (ví dụ: Solar 30 năm vs. Battery 15 năm). Trong đề tài, các thông số CAPEX được tham chiếu từ NREL ATB 2024 \cite{nrel2024atb}.
\fi

% ================================================================
\section{Phân tích kịch bản ngẫu nhiên (Stochastic Scenario Analysis)}
% ================================================================

\subsection{Mô phỏng Monte Carlo}

Mô phỏng Monte Carlo là phương pháp sử dụng lấy mẫu ngẫu nhiên lặp lại để ước lượng các đại lượng thống kê của một hệ thống phức tạp \cite{rubinstein2016monte_carlo}. Thay vì giả định rằng các tham số đầu vào luôn cố định, phương pháp này mô hình hóa chúng như những biến ngẫu nhiên và đánh giá hiệu năng hệ thống trên nhiều lần lấy mẫu khác nhau. Trong bối cảnh Quy hoạch mở rộng nguồn điện (GEP), cách tiếp cận này đặc biệt phù hợp vì bài toán quy hoạch luôn chịu tác động của các yếu tố bất định khó dự báo chính xác trong dài hạn \cite{koltsaklis2018gep,oree2017gep_review}.

Các nguồn bất định quan trọng trong GEP có thể kể đến như sau:
\begin{itemize}
    \item \textbf{Phụ tải điện:} Nhu cầu tiêu thụ thay đổi theo mùa, theo giờ, và chịu ảnh hưởng mạnh của thời tiết, tăng trưởng kinh tế hoặc hành vi người dùng.
    \item \textbf{Nguồn tái tạo:} Điện gió và điện mặt trời phụ thuộc trực tiếp vào điều kiện khí tượng, nên công suất khả dụng có thể dao động lớn ngay cả khi công suất lắp đặt không đổi.
    \item \textbf{Chi phí và thông số công nghệ:} Giá nhiên liệu, hiệu suất vận hành, hệ số công suất, tuổi thọ kinh tế và chi phí bảo trì đều có thể biến thiên theo thời gian.
    \item \textbf{Điều kiện vận hành hệ thống:} Nghịch cảnh truyền tải, sự cố tổ máy hoặc các tình huống thời tiết cực đoan có thể làm suy giảm đáng kể khả năng cung ứng điện.
\end{itemize}

Nếu chỉ tối ưu trên một quỹ đạo dữ liệu duy nhất (deterministic trajectory), phương án đầu tư nhận được có thể đạt chi phí thấp trên bộ dữ liệu đó nhưng lại rất mong manh khi hệ thống phải đối mặt với các trạng thái bất lợi ngoài mẫu. Vì vậy, Monte Carlo được sử dụng để thay thế cách nhìn một-kịch-bản bằng cách đánh giá một phương án đầu tư trên cả một phân phối các trạng thái tương lai có thể xảy ra.

Gọi $\xi$ là véctơ ngẫu nhiên đại diện cho toàn bộ yếu tố bất định của hệ thống, chẳng hạn hồ sơ phụ tải, khả dụng của nguồn tái tạo hoặc các cú sốc ngắn hạn trong vận hành. Khi đó, hàm phạt tổng quát của bài toán không còn là $J(\mathbf{x})$ theo nghĩa xác định, mà trở thành một hàm phụ thuộc vào trạng thái ngẫu nhiên:

\begin{equation}
    J(\mathbf{x}, \xi)
\end{equation}

Mục tiêu của quy hoạch lúc này không còn là tìm nghiệm tốt nhất cho một trạng thái duy nhất, mà là tìm phương án đầu tư có hiệu năng tốt theo một tiêu chí thống kê trên toàn bộ phân phối của $\xi$.

Giá trị kỳ vọng của hàm mục tiêu trên $N$ kịch bản ngẫu nhiên:

\begin{equation}
    \bar{J}(\mathbf{x}) = \frac{1}{N} \sum_{i=1}^{N} J_i(\mathbf{x})
    \label{eq:monte_carlo}
\end{equation}

\noindent trong đó $J_i(\mathbf{x})$ là điểm phạt ứng với kịch bản $i$ của profile phụ tải và năng lượng tái tạo.

Về mặt xác suất, biểu thức trên chính là xấp xỉ mẫu của kỳ vọng toán học:

\begin{equation}
    \mathbb{E}_{\xi}[J(\mathbf{x}, \xi)] \approx \frac{1}{N} \sum_{i=1}^{N} J(\mathbf{x}, \xi_i)
    \label{eq:mc_expectation}
\end{equation}

Trong đó, mỗi $\xi_i$ là một mẫu độc lập hoặc bán độc lập được sinh ra từ cơ chế mô phỏng kịch bản. Khi $N$ tăng lên, trung bình mẫu sẽ hội tụ dần về kỳ vọng thực của hàm mục tiêu theo định luật số lớn \cite{rubinstein2016monte_carlo}. Điều này cho phép biến một bài toán tối ưu dưới bất định thành một bài toán tối ưu trên giá trị trung bình của nhiều lần mô phỏng.

Từ góc nhìn tối ưu hóa, đây là dạng \textit{sample-average approximation}: tầng trên chọn một phương án đầu tư $\mathbf{x}$, tầng dưới đánh giá phương án đó trên nhiều kịch bản ngẫu nhiên, rồi lấy trung bình chi phí để phản hồi lại cho bộ tối ưu. Với cách tổ chức này, Monte Carlo đóng vai trò cầu nối giữa thế giới thực nhiều bất định và mô hình tối ưu hóa cần một hàm đánh giá đủ ổn định để có thể so sánh giữa các phương án đầu tư khác nhau.

Ý nghĩa của phân tích kịch bản ngẫu nhiên không chỉ nằm ở việc làm tăng độ chân thực của bài toán, mà còn ở khả năng bộc lộ rõ các đánh đổi kinh tế -- kỹ thuật. Một phương án đầu tư rất rẻ có thể cho giá trị tốt trong kịch bản cơ sở, nhưng lại gây cắt tải lớn trong các trạng thái hiếm song nghiêm trọng. Ngược lại, một phương án có dự phòng cao hơn có thể làm tăng CAPEX ban đầu nhưng giúp ổn định chi phí vận hành, giảm thiểu rủi ro thiếu điện, và tăng độ bền vững của hệ thống trước các tổ hợp bất lợi của phụ tải và năng lượng tái tạo.

Vì vậy, trong GEP hiện đại, Monte Carlo không đơn thuần là công cụ mô phỏng mà còn là cơ chế định lượng \textit{khả năng chống chịu} của một kế hoạch đầu tư. Nói cách khác, thay vì hỏi ``phương án nào rẻ nhất trong điều kiện lý tưởng?'', phân tích kịch bản ngẫu nhiên đặt ra câu hỏi thực tế hơn: ``phương án nào vẫn vận hành chấp nhận được khi tương lai không diễn ra đúng như kỳ vọng?''.

\subsection{Liên hệ với đề tài}

Trong đề tài này, phân tích kịch bản ngẫu nhiên được sử dụng như nền tảng lý thuyết cho kịch bản \textit{robust} ở các chương sau. Cụ thể, với mỗi véctơ đầu tư $\mathbf{x}$ do tầng trên đề xuất, hệ thống không đánh giá ngay trên một bộ dữ liệu phụ tải -- tái tạo cố định, mà sinh ra nhiều kịch bản Monte Carlo khác nhau để tính:

\begin{equation}
    \bar{J}(\mathbf{x}) = \alpha C_{\text{inv}}(\mathbf{x}) + \frac{1}{N} \sum_{i=1}^{N} \left[ \alpha C_{\text{op}}^*(\mathbf{x}, \xi_i) + \beta S_{\text{total}}(\mathbf{x}, \xi_i) \right]
\end{equation}

Trong đó, phần bất định chủ yếu đến từ nhiễu phụ tải và sự suy giảm công suất khả dụng của điện gió, điện mặt trời. Nhờ đó, tầng trên không còn tối ưu cho một quỹ đạo ``đẹp'' duy nhất mà buộc phải tìm nghiệm có hiệu quả trung bình tốt trên nhiều trạng thái vận hành khác nhau.

Điều này giải thích trực tiếp cho các kết quả thực nghiệm ở chương kết quả: nghiệm robust của đề tài có thể không phải là nghiệm rẻ nhất trên tập ngày đại diện ban đầu, nhưng thường có xu hướng dịch chuyển về những công nghệ điều độ được hơn hoặc có khả năng bù đắp biến thiên tốt hơn, chẳng hạn \texttt{gas\_ccs} và hệ thống lưu trữ. Nói cách khác, Monte Carlo giúp tiêu chí tối ưu phản ánh gần hơn nhu cầu thực của quy hoạch điện: không chỉ tối ưu về mặt chi phí trung bình, mà còn phải đủ bền vững trước các nhiễu động trong tương lai.

% ================================================================
\section{Tổng kết chương}
% ================================================================

Chương này đã trình bày các nền tảng lý thuyết cốt lõi được sử dụng trong đề tài. Bảng~\ref{tab:theory_summary} tóm tắt vai trò của từng lý thuyết trong hệ thống bi-level.

\begin{table}[h!]
\centering
\begin{tabular}{p{4.5cm}p{3cm}p{6cm}}
\toprule
\textbf{Lý thuyết} & \textbf{Vị trí} & \textbf{Vai trò trong đề tài} \\
\midrule
GEP & Toàn bộ hệ thống & Bài toán chính cần giải quyết \\
Bi-level Optimization & Kiến trúc tổng thể & Phân tách quyết định đầu tư và vận hành \\
Linear Programming & Tầng dưới & Giải bài toán điều độ kinh tế \\
Bayesian Optimization & Tầng trên & Tìm phương án đầu tư tối ưu (black-box) \\
MDP & Tầng trên & Khung bài toán cho RL \\
PPO & Tầng trên & Thuật toán RL on-policy ổn định \\
SAC & Tầng trên & Thuật toán RL off-policy với entropy \\
CRF / Niên kim hóa & Hàm mục tiêu & Quy đổi CAPEX thành chi phí hàng năm \\
Monte Carlo & Đánh giá & Phân tích tính bền vững dưới bất định \\
\bottomrule
\end{tabular}
\caption{Tóm tắt vai trò của các lý thuyết trong hệ thống bi-level.}
\label{tab:theory_summary}
\end{table}

Các chương tiếp theo sẽ trình bày chi tiết cách các lý thuyết này được kết hợp và hiện thực hóa trong kiến trúc tối ưu hóa lai hai tầng đề xuất.
